\chapter{Related Work}
\label{ch:related_work}

This chapter reviews prior research on source-code vulnerability detection across static analysis, machine learning methods, and combined workflows. The goal is to position this thesis against existing approaches and identify the specific contribution gap addressed by the proposed system.

\section{Static Analysis for Vulnerability Detection}
\label{sec:rw_static}

Static analysis detects defects by inspecting program structure without running the code. For web-security issues such as SQL injection, XSS, command injection, and CSRF, common techniques include structural pattern matching, AST inspection, and dataflow/taint reasoning.

Rule-based and AST-based tools are widely used in practice. Semgrep~\cite{Semgrep}, for example, matches rules over a program's AST and supports configurable dataflow and taint analysis to trace untrusted inputs to sensitive sinks. Bandit~\cite{Bandit} focuses on Python and checks for common insecure patterns (e.g., use of \texttt{eval}, hardcoded passwords). Pylint and similar linters can be extended with security-oriented rules. These tools are fast and interpretable but depend on the quality and coverage of rules; they may miss novel variants or produce false positives when patterns are too broad. In industry, additional static analysis platforms such as CodeQL~\cite{GitHubCodeQL}, SonarQube~\cite{SonarQube}, and Coverity~\cite{Coverity} are commonly used to detect security issues and code quality problems across large codebases.

Studies on developer adoption show that static analysis tools are underused partly due to usability and trust: developers want actionable, relevant results and are frustrated by noisy or hard-to-interpret output~\cite{Christakis2016,Johnson2013}. This motivates tools that combine multiple analysis strategies and present findings in a clear, integrated way, as in the static analysis component of this thesis.

\section{Machine Learning for Code Vulnerability Detection}
\label{sec:rw_ml}

Machine learning approaches infer vulnerability patterns from large code corpora and historical commits. Recent work combines learned code representations with deep sequence or graph models for vulnerability classification~\cite{Allamanis2018Survey,Alon2019code2vec,Feng2020CodeBERT,Zhou2019Devign,SySeVR2021,Steenhoek2024DeepDFA}.

Farasat and Posegga~\cite{Farasat2024MLPython} present a comparative study of machine learning techniques for Python source code vulnerability detection (CODASPY~'24). They compare Gaussian Naive Bayes (GNB), Decision Tree, Logistic Regression (LR), Multi-Layer Perceptron (MLP), and BiLSTM. Their results show that BiLSTM combined with Word2Vec achieves state-of-the-art performance, with average accuracy of approximately 98.6\% and F1-score of approximately 94.7\%. The open-source code and model architecture are available in the Python-Source-Code-Vulnerability-Detection repository~\cite{PythonVulnDetectionDataset}, which uses a TensorFlow-based BiLSTM setup for the same vulnerability types (SQL injection, XSS, command injection, CSRF).

Recurrent architectures such as LSTM and BiLSTM are commonly used for sequence classification and can be trained to predict whether a code snippet is vulnerable. The VulnerabilityDetection-master project~\cite{VulnDetection} and the Python-Source-Code-Vulnerability-Detection repository~\cite{PythonVulnDetectionDataset} provide datasets and pretrained BiLSTM models for Python vulnerability types (SQL injection, XSS, command injection, CSRF), mined from real-world GitHub commits. Other work explores datasets such as DiverseVul for C/C++ vulnerabilities and evaluates a range of neural architectures~\cite{Fan2023DiverseVul,Zhou2019Devign}. Word2Vec and similar embeddings~\cite{Word2VecRef}, as well as more recent code representation models like code2vec and CodeBERT~\cite{Alon2019code2vec,Feng2020CodeBERT}, are often used to obtain distributed representations of code tokens or AST paths, which are then fed into LSTM-based or graph-based models for classification or localization~\cite{LSTMRef,Li2018VulDeePecker,Zhang2025DLSurvey}. Such models can generalize to unseen variants that follow similar patterns but require substantial training data and may lack interpretability compared to rule-based analysis.

\subsection{Other Python-Focused ML Approaches}
\label{sec:rw_ml_python}

Additional studies focus directly on Python vulnerability detection. Wartschinski et al.~\cite{VulnDetection} present VUDENC, which applies machine learning to Python code mined from natural repositories. Wang et al.~\cite{Wang2024DLPythonVuln} compare multiple deep architectures for Python vulnerability detection, and Bagheri and Hegedűs~\cite{Bagheri2021PythonVuln} evaluate code-representation choices using BERT embeddings with LSTM-based sequence modeling.

Table~\ref{tab:rw_python_ml} summarizes these Python-focused ML approaches by key idea or model.

\begin{table}[htbp]
\centering
\caption{Selected Python-focused ML approaches for vulnerability detection.}
\label{tab:rw_python_ml}
\begin{tabular}{@{}ll@{}}
\hline
\textbf{Approach} & \textbf{Key idea / model} \\
\hline
Wartschinski et al. & VUDENC: ML-based Python vulnerability detection \\
Wang et al. & Deep learning / NLP models for Python code \\
Bagheri \& Hegedűs & BERT + LSTM for sequence/code analysis \\
\hline
\end{tabular}
\end{table}

\section{Comparison and Hybrid Approaches}
\label{sec:rw_comparison}

Only a limited number of studies compare static analysis and machine learning on the same vulnerability categories within one user-facing platform. Static analysis typically provides immediate and explainable findings without training data, while ML can capture subtler patterns but introduces data dependence and reduced interpretability. Dual-mode systems are therefore useful both in practice and in research because they enable side-by-side comparison of these trade-offs.

The system in this thesis provides exactly such a comparison: handcrafted static analysis (regex, AST, taint flow) and ML-based scoring using pretrained BiLSTM models, with a unified interface for file upload, paste, or GitHub URL. The evaluation in the main body compares precision, recall, and F1 across vulnerability types and discusses trade-offs and ways to reduce false positives and improve interpretability.
